To develop, test, and deploy the 3D Lab Manager successfully, the team requires regular access to the senior design labs. The primary physical workspace will be ERB202. Because the application is inherently designed to manage this specific space, the team will need access to map the physical layout, track current inventory locations, and understand the practical bottlenecks the lab manager faces on a daily basis. The other labs will also be used for mapping specific storage sites of lab equipment as well. 

From an infrastructure perspective, the project will be hosted on the university's senior design network. Steven, the lab manager, will assign a dedicated Linux box to the team. This machine will serve as the localized production environment housing the core application server and the redundant local database. The team will require SSH access and administrative privileges on this Linux box to configure the disaster recovery protocols, setup the routing architecture, and deploy the backend logic. We will also be utilizing a 3D scanner to get the necessary imaging of the labs.

For software needs, we will utilize the legacy 3D image files and the previous GitHub repository to establish a baseline. However, to generate and modify the 3D models of the listed labs, we will be installing and utilizing blender so that we can improve on the models that were received from the previous groups'  iteration.

For the kiosk, we will need to buy a touch-screen device so that users can interact with it to find out the information on where things exist within the labs.
